\section{Applications}
\label{sec-applications}

\subsection{Intrinsic Use: Using \eclector as the Reader of a \cl{} Implementation}

The most obvious application for \eclector{} is as a native reader of
some \cl{} implementation.  Currently, \eclector{} is used this way
only in Clasp, but it will also be used in the SICL \cl{}
implementation, once it is finished.

Currently, SICL uses \eclector{} extrinsically during bootstrapping,
in order to read forms to be evaluated in a first-class global
environment \cite{Strandh:2015:ELS:Environments}.  Since the host
global environment must not be affected as a result of reading these
forms, we take full advantage of the ability to configure \eclector{}
to avoid such side effects.

\subsection{Alternate Syntax beyond Readtable and Reader Macros}

https://github.com/pve1/eclector-access/

maybe coalton

maybe Hayley's language

\subsection{Reader Part of Sandboxed Evaluation}

\subsection{Syntax Highlighting}

For documentation, editors, REPLs and IDEs

Important features

\begin{itemize}
\item No side effects (Interning, package system changes, code evaluation)
\item Error recovery
\end{itemize}

\subsection{Static Analysis}

Important features

\begin{itemize}
\item No side effects (Interning, package system changes, code evaluation)
\item Error recovery
\end{itemize}

Limitations

\begin{itemize}
\item Reader macros
\item Read-time evaluation
\end{itemize}

TODO link to fukamachi's linter, atgreen's linter
